\documentclass[11pt]{article}
\usepackage[T1]{fontenc}
\usepackage[utf8]{inputenc}
\usepackage{lmodern,microtype}
\usepackage[a4paper,margin=27mm]{geometry}
\usepackage{amsmath,amssymb,amsthm,mathtools,booktabs,array,xcolor}
\usepackage[numbers,sort&compress]{natbib}
\usepackage{hyperref}
\hypersetup{colorlinks=true,linkcolor=blue!45!black,citecolor=green!35!black,
urlcolor=blue!55!black,pdfauthor={},
pdftitle={Complexity-Weighted Moments, Critical Mass Escape, and Selected Corona States}}
\newtheorem{theorem}{Theorem}[section]
\newtheorem{proposition}[theorem]{Proposition}
\newtheorem{lemma}[theorem]{Lemma}
\newtheorem{corollary}[theorem]{Corollary}
\theoremstyle{definition}
\newtheorem{definition}[theorem]{Definition}
\newtheorem{example}[theorem]{Example}
\newcommand{\N}{\mathbb N_0}
\newcommand{\C}{\mathbb C}
\newcommand{\R}{\mathbb R}
\newcommand{\E}{\mathbb E}
\newcommand{\A}{\mathcal A}
\newcommand{\K}{\mathcal K}
\newcommand{\eps}{\varepsilon}
\newcommand{\norm}[1]{\left\lVert#1\right\rVert}
\newcommand{\ind}{\mathbf 1}
\DeclareMathOperator{\Res}{Res}
\title{\textbf{Complexity-Weighted Moments, Critical Mass Escape,\\
and Selected Corona States}}
\author{}
\date{29 October 2026}
\begin{document}
\maketitle
\begin{abstract}
A positive sum over graded sectors can have determinate moments at every fixed regulator and still lose all finite-level probability at its critical boundary.
We give a self-contained mathematical methods treatment of this distinction.
Exponential counting and intensity bounds imply a finite exponential moment, while a heat-damped principal-series example admits an exact normalizer and a support-safe proof for translated clocks.
At the true convergence abscissa, finite critical normalization gives total-variation convergence; divergent normalization is equivalent to escape from every bounded complexity set.
Escape forces divergence of all positive complexity moments, although rates require additional hypotheses.
Two-sided exponential level bounds give order estimates, and a convergent residual profile gives an exponential scaling law.
Bounded sector states descend to a vanishing-ideal quotient, but word and combined word/microlocal quotients need not agree.
We prove a corrected further-quotient construction for word-escaped states.
Finally, geometric Abel regulators fix parity at one half and exclude ultrafilter characters, whereas explicit long-block observables give two genuinely reachable separating cluster states.
The contribution is an integrated analysis and corrected set of examples using classical moment, summability and operator-algebraic ingredients.
It does not assert a new general determinacy theorem or construct a physical gravitational completion.
\end{abstract}

\section{The question and its classical ingredients}

Three uniqueness questions arise naturally for a regulated positive ensemble.
Does the entire moment sequence of a nonnegative clock determine its scalar probability law?
Does the family of normalized laws have a limit when the regulator reaches its convergence boundary?
If limiting states exist on a larger algebra of bounded observables, do the selected states agree on the observables of interest?
The answers concern different mathematical objects.
Neither a divergent moment nor a large ambient state space settles the other questions.

Here a complexity function is a grading or exhaustion, rather than computational, circuit or holographic complexity.
For nonnegative intensities $a_\gamma$ and integer grading $\ell(\gamma)$, the probabilities are
\begin{equation}
 p_t(\gamma)=\frac{e^{-t\ell(\gamma)}a_\gamma}{Z(t)},
 \qquad Z(t)=\sum_\gamma e^{-t\ell(\gamma)}a_\gamma.
 \label{eq:sector-start}
\end{equation}
At fixed sufficiently large $t$, the clock $\ell$ has a finite exponential moment.
Near criticality, normalization may remain finite or diverge.
The latter case produces weak-* states at infinity, depending on the observable algebra and on which states the regulator actually reaches.

The ingredients have established antecedents.
Carleman's criterion and exponential-moment determinacy belong to classical moment theory; Simon distinguishes the Stieltjes support problem from the Hamburger problem~\cite{Simon1998}.
Patterson's normalized critical series provide a geometric instance of a divergent denominator concentrating limit measures on a boundary~\cite{Patterson1976}.
Banach limits and almost convergence, including the unique parity value under shift-invariant positive limits, were studied by Lorentz~\cite{Lorentz1948}.
Weak-* compactness and quotient-state facts are standard functional analysis~\cite{Alaoglu1940,Murphy1990}.
The proofs below specialize and combine these ideas; we make no claim that the general ingredients originate here.

The combined treatment concerns the passage from a fixed-regulator moment problem to a selected boundary-state problem.
Two local distinctions are consequential.
A positive translation of a Stieltjes clock cannot be justified by an unrestricted inverse translation of competing measures.
A scalar function vanishing at word infinity need not vanish along a microlocal direction with the word fixed.
We resolve both issues explicitly.

The heat example uses a spectral density from Virasoro topological quantum field theory~\cite{CollierEberhardtZhang2023}.
Only the positive one-variable density is imported.
No physical state selection or gravitational interpretation is inferred from it.
Thermal point-process statistics are a separate question: the objects here are positive graded sums, critical normalization and corona limits.

\section{Determinacy at a fixed regulator}
\label{sec:determinacy}

\subsection{Moment conventions and the exponential criterion}

\begin{definition}
Let $\mu$ be a probability on $[0,\infty)$ with finite moments $m_n=\int x^n\,d\mu(x)$ for all $n\in\N$.
Its Stieltjes moment problem is determinate if no other positive Borel measure on $[0,\infty)$ has the same entire moment sequence.
The Hamburger problem permits competitors on all of $\R$.
\end{definition}

We use the Stieltjes--Carleman sufficient condition
\begin{equation}
 \sum_{n=1}^\infty m_n^{-1/(2n)}=\infty.
 \label{eq:carleman}
\end{equation}
A vanishing positive moment forces the measure to be concentrated at zero, so that degenerate case is immediate.
Failure of~\eqref{eq:carleman} does not establish indeterminacy.
The classical criterion is reviewed in~\cite[Corollary 4.5]{Simon1998}.

\begin{lemma}[Finite exponential moment]
\label{lem:mgf}
If $X\geq0$ and $M_s=\E e^{sX}<\infty$ for some $s>0$, all its moments are finite and its Stieltjes moment problem is determinate.
\end{lemma}
\begin{proof}
The nonnegative exponential series gives
\begin{equation}
 x^n\leq n!s^{-n}e^{sx},\qquad x\geq0.
\end{equation}
Thus $m_n\leq n!s^{-n}M_s$.
Stirling's formula implies
\begin{equation}
 m_n^{-1/(2n)}\geq
 s^{1/2}(n!)^{-1/(2n)}M_s^{-1/(2n)}
 =s^{1/2}e^{1/2}n^{-1/2}(1+o(1)).
\end{equation}
The comparison series diverges, so~\eqref{eq:carleman} applies.
\end{proof}

This proves uniqueness of the scalar law of $X$.
States on a larger noncommutative algebra can share that law and differ elsewhere.
No informational completeness of the entire state is asserted.

\subsection{A heat-damped spectral example and translated clocks}

For $b>0$ and $P\geq0$, set
\begin{equation}
 \rho_b(P)=4\sqrt2\sinh(2\pi bP)\sinh(2\pi b^{-1}P),
 \qquad Q=b+b^{-1}.
 \label{eq:density}
\end{equation}
This is the convention of~\cite[Eq.~(2.23)]{CollierEberhardtZhang2023}.
For $\beta>0$, define
\begin{equation}
 Z_\beta=\int_0^\infty e^{-\beta P^2}\rho_b(P)\,dP,
 \qquad d\mu_\beta=Z_\beta^{-1}e^{-\beta P^2}\rho_b(P)\,dP.
 \label{eq:heat}
\end{equation}

\begin{proposition}[Heat normalization and determinacy]
\label{prop:heat}
For $b,\beta>0$, $0<Z_\beta<\infty$, with exact value
\begin{equation}
 Z_\beta=\sqrt{\frac{2\pi}{\beta}}
 \left[e^{\pi^2(b+b^{-1})^2/\beta}
       -e^{\pi^2(b-b^{-1})^2/\beta}\right].
 \label{eq:normalizer}
\end{equation}
Under $\mu_\beta$, each of $P^2$, $P^2+c$ for fixed $c\geq0$, and $\kappa P$ for fixed $\kappa>0$ is Stieltjes determinate.
For $X=P^2$ and $0<s<\beta$,
\begin{equation}
 \E_\beta e^{s(X+c)}=e^{sc}\frac{Z_{\beta-s}}{Z_\beta}.
 \label{eq:translated-mgf}
\end{equation}
\end{proposition}
\begin{proof}
Since $\sinh u\leq e^u/2$ for $u\geq0$, one has
$0\leq\rho_b(P)\leq\sqrt2 e^{2\pi QP}$.
Completing the square bounds the heat integral by a finite Gaussian integral; positivity holds since the integrand is positive for $P>0$.
For the exact formula use
\begin{equation}
 2\sinh u\sinh v=\cosh(u+v)-\cosh(u-v),\qquad
 \int_0^\infty e^{-\beta P^2}\cosh(aP)\,dP
 =\frac{\sqrt\pi}{2\sqrt\beta}e^{a^2/(4\beta)}.
\end{equation}
Substitution gives~\eqref{eq:normalizer}.
The bracket is strictly positive since
$(b+b^{-1})^2-(b-b^{-1})^2=4$.
Equation~\eqref{eq:translated-mgf} follows directly and is finite for $s<\beta$.
Lemma~\ref{lem:mgf} proves determinacy of $X$ and $X+c$.
For $L=\kappa P$, all finite positive exponential moments exist since
$-\beta P^2+(2\pi Q+s\kappa)P$ is Gaussian dominated.
The same lemma proves the length-clock claim.
\end{proof}

The translation step requires care.
An arbitrary competitor $\nu$ for $Y=X+c$ on $[0,\infty)$ may translate into $[-c,\infty)$, with mass below zero.
Stieltjes uniqueness for $X$ alone does not exclude it.
General Stieltjes determinacy is not translation invariant; the distinction is explicit in~\cite[Section 3 and Appendix C]{Simon1998}.
Equation~\eqref{eq:translated-mgf} proves the actual heat-family conclusion without that support gap.

Gaussian domination also gives, for a finite constant $C_\beta$ independent of $n$,
\begin{equation}
 \E_\beta P^{2n}\leq C_\beta
 \left(\frac{2}{\beta}\right)^{n+1/2}
 \Gamma\!\left(n+\frac12\right).
\end{equation}
Indeed $-\beta P^2+2\pi QP\leq-(\beta/2)P^2+(2\pi Q)^2/(2\beta)$, and the remaining integral is evaluated by $u=\beta P^2/2$.
This supplies the same Carleman divergence.
The bare density is not a probability, and no conclusion removes the regulator or selects a physical $\beta$.

\subsection{Exponentially bounded sector ensembles}

Let $\Gamma$ be countable, $\ell:\Gamma\to\N$, and
$N_R=\#\{\gamma:\ell(\gamma)=R\}$.
Assume finite nonnegative intensities, not all zero, and constants $A,C<\infty$, $h,r\geq0$ such that
\begin{equation}
 N_R\leq Ce^{hR},\qquad a_\gamma\leq Ae^{r\ell(\gamma)}.
 \label{eq:shell-bounds}
\end{equation}

\begin{theorem}[Supercritical determinacy]
\label{thm:supercritical}
Under~\eqref{eq:shell-bounds},~\eqref{eq:sector-start} defines a probability for every $t>h+r$.
The clock $L=\ell$ has a finite exponential moment and is Stieltjes determinate.
\end{theorem}
\begin{proof}
Grouping by level gives
\begin{equation}
 0<Z(t)\leq CA\sum_{R\geq0}e^{-(t-h-r)R}<\infty.
\end{equation}
For $0<s<t-h-r$,
\begin{equation}
 \E_t e^{sL}\leq\frac{CA}{Z(t)}
 \sum_{R\geq0}e^{-(t-h-r-s)R}<\infty.
\end{equation}
Apply Lemma~\ref{lem:mgf}.
\end{proof}

If $a_\gamma=\norm{V_\gamma\eta}^2$ and
$\norm{V_\gamma\eta}\leq C_\eta e^{r_0\ell(\gamma)}$, the intensity exponent is $r=2r_0$.
An action satisfying $A_{\mathrm{act}}(\gamma)\geq t\ell(\gamma)-B$ gives the same bound up to $e^B$.
These are sufficient comparisons, not identities of the action with the grading.
In particular, $h+r$ need not be the true critical exponent.

\section{Critical normalization, escape and moment rates}
\label{sec:critical}

Let $B_R\geq0$ be finite, with infinitely many nonzero levels, and define
\begin{equation}
 Z(t)=\sum_{R\geq0}e^{-tR}B_R,
 \qquad
 \alpha=\limsup_{\substack{R\to\infty\\B_R>0}}
 \frac{\log B_R}{R}\in\R.
 \label{eq:abscissa}
\end{equation}
The finite-real-$\alpha$ assumption excludes superexponential growth and the case $\alpha=-\infty$ from this critical-boundary analysis.
Whenever $Z(t)<\infty$, put
\begin{equation}
 q_t(R)=Z(t)^{-1}e^{-tR}B_R.
 \label{eq:level-law}
\end{equation}

\begin{lemma}[The true abscissa]
\label{lem:abscissa}
$Z(t)<\infty$ for $t>\alpha$, and $Z(t)=\infty$ for $t<\alpha$.
At $t=\alpha$ either behavior is possible.
Every fixed $t>\alpha$ has a determinate complexity law.
\end{lemma}
\begin{proof}
For $t>\alpha$, choose $\delta>0$ with $\alpha+\delta<t$.
Eventually $B_R\leq e^{(\alpha+\delta)R}$, giving geometric domination.
For $t<\alpha$, infinitely many $R$ satisfy $B_R\geq e^{(t+\delta)R}$ for a fixed $\delta>0$.
Those summands do not tend to zero.
The profiles $e^{\alpha R}(R+1)^{-2}$ and $e^{\alpha R}(R+1)^{-1}$ give the two critical behaviors.
Finally, $0<s<t-\alpha$ gives
$\E_t e^{sL}=Z(t-s)/Z(t)<\infty$, so Lemma~\ref{lem:mgf} applies.
\end{proof}

\begin{definition}
A net of probabilities $q_i$ on $\N$ has mass escape if
\begin{equation}
 q_i(\{0,\ldots,K\})\longrightarrow0
 \quad\text{for every finite }K.
 \label{eq:escape}
\end{equation}
\end{definition}

\begin{theorem}[Critical-denominator alternative]
\label{thm:alternative}
Precisely one branch holds as $t\downarrow\alpha$:
\begin{enumerate}
\item If $Z(\alpha)<\infty$, then $q_t$ tends in total variation to
$q_\alpha(R)=Z(\alpha)^{-1}e^{-\alpha R}B_R$.
\item If $Z(\alpha)=\infty$, then $q_t$ has mass escape and has no weak probability limit on the original discrete space $\N$.
\end{enumerate}
\end{theorem}
\begin{proof}
In the first branch put $w_t(R)=e^{-tR}B_R$.
Dominated convergence gives $\norm{w_t-w_\alpha}_1\to0$ and $Z(t)\to Z(\alpha)>0$.
Then
\begin{equation}
 \norm{q_t-q_\alpha}_1\leq
 \frac{\norm{w_t-w_\alpha}_1}{Z(t)}
 +Z(\alpha)\left|\frac1{Z(t)}-\frac1{Z(\alpha)}\right|\longrightarrow0.
\end{equation}
Our convention is $d_{\mathrm{TV}}(p,q)=\frac12\norm{p-q}_1$.
In the second branch monotone convergence gives $Z(t)\uparrow\infty$.
For fixed $K$, the numerator $\sum_{R\leq K}e^{-tR}B_R$ stays bounded, so its ratio to $Z(t)$ vanishes.
A weak probability limit on the discrete space would give every singleton mass zero, contradicting countable additivity and normalization.
\end{proof}

\begin{proposition}[Minimal escape criterion]
\label{prop:minimal}
Let $t_c\in\R$, every $B_R$ be finite, and $0<Z(t)<\infty$ for $t>t_c$.
Then $Z(t)\to\infty$ as $t\downarrow t_c$ if and only if $q_t$ has mass escape.
\end{proposition}
\begin{proof}
The forward implication is the bounded finite-level numerator argument.
For the converse choose $R_0$ with $B_{R_0}>0$, which exists since $Z(t)>0$.
A bounded subnet of $Z(t)$ as $t\downarrow t_c$ would keep
$q_t(R_0)=e^{-tR_0}B_{R_0}/Z(t)$ bounded away from zero, contradicting escape.
\end{proof}

The denominator condition covers sparse and irregular residual weights, if the actual critical series diverges.
The basic boundary-support mechanism also appears in~\cite[Section 3]{Patterson1976}; its geometric and conformal conclusions are not assumed here.

\begin{proposition}[Every positive moment diverges under escape]
\label{prop:moment-divergence}
If~\eqref{eq:escape} holds, then $\E_i L^p\to\infty$ for every $p>0$, with extended-real expectations allowed.
\end{proposition}
\begin{proof}
For every $K>0$, $\E_i L^p\geq K^p q_i(L\geq K)$.
The probability tends to one, so $\liminf_i\E_i L^p\geq K^p$ for every $K$.
Let $K\to\infty$.
\end{proof}

Tail regularity is needed for a rate, not for this divergence.
Conversely, the finite-denominator profile $(R+1)^{-2}$ defines an ordinary critical probability while its positive integer moments diverge.
Moment divergence and escape are inequivalent.

\begin{theorem}[Order estimates and a regular-profile limit]
\label{thm:rates}
Suppose for all sufficiently large $R$,
\begin{equation}
 0<c_-\leq c_R:=e^{-\alpha R}B_R\leq c_+<\infty.
 \label{eq:two-sided}
\end{equation}
Writing $t=\alpha+\eps$, one has
\begin{equation}
 Z(\alpha+\eps)=\Theta(\eps^{-1}),\qquad
 \E_\eps L^n=\Theta(\eps^{-n})
 \quad(n\geq1\text{ fixed}).
 \label{eq:theta}
\end{equation}
If in addition $c_R\to c\in(0,\infty)$, then
\begin{equation}
 Z(\alpha+\eps)\sim\frac c\eps,
 \qquad \eps^n\E_\eps L^n\longrightarrow n!,
 \label{eq:asymptotic}
\end{equation}
and $\eps L$ converges in law to a mean-one exponential variable.
\end{theorem}
\begin{proof}
The finite initial part is negligible against
$\sum e^{-\eps R}=(1-e^{-\eps})^{-1}=\Theta(\eps^{-1})$.
The moment numerator is compared above and below with $\sum R^n e^{-\eps R}$.
The Riemann-sum limit
\begin{equation}
 \eps^{n+1}\sum_{R\geq0}R^n e^{-\eps R}
 =\eps\sum_{R\geq0}(\eps R)^n e^{-\eps R}
 \longrightarrow\int_0^\infty u^n e^{-u}\,du=n!
 \label{eq:riemann}
\end{equation}
then proves~\eqref{eq:theta} after normalization.
To justify the unbounded interval, first use Riemann sums on a compact interval.
The remaining tail is uniformly bounded by a constant multiple of the integrable tail of $(1+u)^n e^{-u}$.
If $c_R\to c$, split at fixed large $K$.
The finite part vanishes after multiplication by $\eps^{n+1}$.
Above $K$, replacement of $c_R$ by $c$ has error bounded by
$\sup_{R>K}|c_R-c|$ times the left side of~\eqref{eq:riemann}.
Let $\eps\downarrow0$ and then $K\to\infty$, obtaining~\eqref{eq:asymptotic}, including $n=0$ for normalization.
For $x\geq0$, the same tail Riemann sum gives
\begin{equation}
 \Pr_\eps(\eps L>x)=
 \frac{\sum_{R>x/\eps}c_R e^{-\eps R}}
 {\sum_{R\geq0}c_R e^{-\eps R}}
 \longrightarrow\int_x^\infty e^{-u}\,du=e^{-x}.
\end{equation}
This proves the distributional limit.
\end{proof}

Section~\ref{sec:abel} gives a two-sided bounded residual profile with two different subsequential values of $\eps Z$.
Thus~\eqref{eq:theta} cannot generally be upgraded to a constant equivalent.
Every fixed regulated law remains determinate by Lemma~\ref{lem:abscissa} even as all positive moments diverge at the critical boundary.

\section{Bounded observable algebras and corona limits}
\label{sec:corona}

Write $\ell^\infty=\ell^\infty(\N)$ for bounded complex sequences, $c_0$ for sequences tending to zero, and $c$ for convergent sequences.
The scalar corona is the unital $C^*$-algebra $\ell^\infty/c_0$.
A state is a positive complex-linear functional $\omega$ with $\omega(1)=1$; it obeys $|\omega(a)|\leq\norm a$.
Weak-* convergence means convergence on every fixed algebra element.
State spaces are weak-* compact by Banach--Alaoglu and weak-* closedness of positivity and normalization~\cite{Alaoglu1940,Murphy1990}.
For nonseparable algebras this supplies subnets, not necessarily convergent subsequences.
Each probability $q_i$ defines a state
\begin{equation}
 \omega_i(x)=\sum_{R\geq0}q_i(R)x_R,\qquad x\in\ell^\infty.
\end{equation}

\begin{theorem}[Tail values and scalar descent]
\label{thm:tail}
Under mass escape, $\omega_i(f)\to0$ for $f\in c_0$, and every weak-* cluster state factors through $\ell^\infty/c_0$.
If $x_R\to L$, then $\omega_i(x)\to L$.
All cluster states therefore have the same restriction to $c$, and $c/c_0\simeq\C$ has one state.
\end{theorem}
\begin{proof}
Given $\eta>0$, choose $K$ with $|x_R-L|<\eta$ for $R>K$.
Then
\begin{equation}
 |\omega_i(x)-L|\leq\norm{x-L}_\infty q_i(R\leq K)+\eta.
\end{equation}
First take the escape limit, then $\eta\downarrow0$.
Taking $L=0$ gives annihilation of $c_0$.
A functional annihilating a closed ideal defines a unique functional on the quotient; positivity and unitality pass to that functional.
The limit map $c\to\C$ is a surjective unital $*$-homomorphism with kernel $c_0$, proving the final assertion.
\end{proof}

\begin{proposition}[Ambient tail separation]
\label{prop:ultrafilter}
A real nonconvergent $x\in\ell^\infty$ is separated by two states on $\ell^\infty/c_0$ that agree on every convergent tail.
\end{proposition}
\begin{proof}
Choose subsequences tending to distinct cluster values $a<b$ and free ultrafilters containing their index sets.
The ultralimits $\phi_\mathcal U(y)=\lim_\mathcal U y_R$ are positive, linear and unital, annihilate $c_0$, and give $x$ the values $a,b$.
Both agree with the ordinary limit on convergent sequences.
\end{proof}

Free ultrafilters use a corresponding choice principle.
This argument identifies ambient states, without showing that the regulator selects either one.
Any unital scalar subalgebra contained in $c/c_0\simeq\C$ has no state-separating observable.
This conclusion does not exclude separation on other ambient noncommutative observables.

\subsection{Sector products}

Let $\A_\gamma$ be nonzero unital $C^*$-algebras over a countable graded set.
Define the bounded product and length-vanishing ideal by
\begin{align}
 \A&=\{a=(a_\gamma):a_\gamma\in\A_\gamma,
                     \ \sup_\gamma\norm{a_\gamma}<\infty\},\\
 I_\ell&=\{a\in\A:
       \sup_{\ell(\gamma)>K}\norm{a_\gamma}\to0
       \text{ as }K\to\infty\}.
 \label{eq:length-ideal}
\end{align}
Operations are coordinatewise.
This is a closed self-adjoint two-sided ideal: uniform convergence preserves the tail condition and multiplication by bounded families preserves vanishing.
If bounded length balls are finite, $I_\ell$ is the usual $c_0$ direct-sum ideal.
Otherwise it is a length-based ideal, not silently ordinary coordinatewise $c_0$.

\begin{proposition}[Sector-state descent]
\label{prop:sector}
Let $p_i$ be probabilities with $p_i(\ell\leq K)\to0$ for every $K$, and let $\sigma_\gamma$ be sector states.
Then
\begin{equation}
 \Omega_i(a)=\sum_\gamma p_i(\gamma)\sigma_\gamma(a_\gamma)
\end{equation}
is a state on $\A$, and every weak-* cluster state annihilates $I_\ell$ and descends to $\A/I_\ell$.
\end{proposition}
\begin{proof}
Absolute convergence follows from $|\sigma_\gamma(a_\gamma)|\leq\norm a$.
Positivity and normalization hold term by term.
For $a\in I_\ell$, choose $K$ with $\sup_{\ell(\gamma)>K}\norm{a_\gamma}<\eta$.
Then $|\Omega_i(a)|\leq\norm a\,p_i(\ell\leq K)+\eta$.
Take the escape limit and then $\eta\downarrow0$.
The quotient conclusion follows as in Theorem~\ref{thm:tail}.
\end{proof}

The bound applies entrywise to finite matrix amplifications.
It does not prove that an independently specified map $\Phi$ descends: that requires $\Phi(I_\ell)\subseteq I_\ell$.
Complete positivity of a descended map requires complete positivity of the map being descended.
Unbounded observables require separately declared domains and graph estimates.
State descent, operation descent and unbounded-domain control are distinct assertions.

\section{Word and combined escape: a corrected quotient}
\label{sec:word}

Let $j=(\gamma,k)$, with each countable fiber $K_\gamma$ nonempty.
Assume $\ell_W:\Gamma\to\N$ has finite bounded balls and
\begin{equation}
 \ell_T(\gamma,k)\geq\ell_W(\gamma).
 \label{eq:projection}
\end{equation}
For example, total complexity may be word plus microlocal complexity.
No physical meaning of $k$ is required.
On the bounded product $\A=\prod_{\gamma,k}\A_{\gamma,k}$ of nonzero unital factors, define
\begin{align}
 I_T&=\{a:\sup_{\ell_T(\gamma,k)>R}\norm{a_{\gamma,k}}\to0\},\\
 J_W&=\{a:\sup_{\ell_W(\gamma)>R,\ k\in K_\gamma}
                         \norm{a_{\gamma,k}}\to0\}.
 \label{eq:two-ideals}
\end{align}
Both are closed two-sided ideals.

\begin{example}[A word-vanishing lift can survive total escape]
\label{ex:ray}
Take scalar factors on $\Gamma=\N$, $K_\gamma=\N$, with
$\ell_W(\gamma)=\gamma$ and $\ell_T(\gamma,k)=\gamma+k$.
Although $f(\gamma)=\ind_{\{0\}}(\gamma)$ belongs to word $c_0$, its lift equals one on $(0,k)$ for every $k$.
Thus $\widetilde f\notin I_T$, and scalar pullback does not generally define a map from the word corona to $\A/I_T$.
\end{example}

Inequality~\eqref{eq:projection} says a bounded total ball projects into a bounded word ball.
It does not bound the microlocal fibers over a fixed word.
These directions must not be interchanged.

\begin{theorem}[Further quotient for word-escaped states]
\label{thm:word-quotient}
Under~\eqref{eq:projection}, $I_T\subseteq J_W$.
Suppose the normalized states are
\begin{equation}
 \Omega_i(a)=\sum_\gamma p_i(\gamma)
 \sum_{k\in K_\gamma}\nu_i(k\mid\gamma)
                 \sigma_{\gamma,k}(a_{\gamma,k}),
 \label{eq:conditional-lift}
\end{equation}
where $p_i$ are word-escaped probabilities, $\nu_i(\cdot\mid\gamma)$ are conditional probabilities and $\sigma_{\gamma,k}$ are sector states.
Every cluster state annihilates $J_W$ and factors through $\A/J_W$, a further quotient of $\A/I_T$.
Scalar lifting induces a central isometric unital embedding
\begin{equation}
 \ell^\infty(\Gamma)/c_0(\Gamma)\lhook\joinrel\longrightarrow\A/J_W.
 \label{eq:word-embedding}
\end{equation}
\end{theorem}
\begin{proof}
Word-large implies total-large by~\eqref{eq:projection}, so the word-tail supremum is bounded by the total-tail supremum.
Thus $I_T\subseteq J_W$.
For $a\in J_W$, splitting~\eqref{eq:conditional-lift} gives
\begin{equation}
 |\Omega_i(a)|\leq\norm a\,p_i(\ell_W\leq R)
       +\sup_{\ell_W(\gamma)>R,k}\norm{a_{\gamma,k}}.
\end{equation}
Take the regulator limit and then $R\to\infty$ to obtain annihilation.

Lift $f$ by $\widetilde f_{\gamma,k}=f(\gamma)1_{\gamma,k}$, a central unital $*$-homomorphism.
Properness of $\ell_W$ makes a scalar word-$c_0$ lift belong to $J_W$, so it descends.
Its quotient norm is
\begin{equation}
 \norm{\widetilde f+J_W}
 =\lim_{R\to\infty}\sup_{\ell_W(\gamma)>R}|f(\gamma)|
 =\norm{f+c_0(\Gamma)}.
 \label{eq:quotient-norm}
\end{equation}
For the upper bound cancel the lift on $\ell_W\leq R$ by an element of $J_W$.
For the lower bound, every $a\in J_W$ satisfies
\begin{equation}
 \norm{\widetilde f+a}\geq
 \sup_{\ell_W(\gamma)>R}|f(\gamma)|
 -\sup_{\ell_W(\gamma)>R,k}\norm{a_{\gamma,k}}.
\end{equation}
Let $R\to\infty$ and infimize over $a$.
This proves isometry and injectivity; centrality and unitality pass to the quotient.
\end{proof}

In Example~\ref{ex:ray}, $I_T\subsetneq J_W$.
The fixed-word lift survives modulo $I_T$, but every word-escaped state assigns it zero.
The further quotient records that additional annihilation.
If word balls are not finite, replace $c_0(\Gamma)$ in~\eqref{eq:word-embedding} by the explicitly length-vanishing scalar ideal.

A stronger geometric condition restores scalar pullback into the total quotient: every fiber over a bounded word ball must have bounded total complexity.
Then, for $f$ word-$c_0$, choose a word ball outside which $|f|<\eta$.
Beyond the total bound of its preimage, all remaining lifted values are smaller than $\eta$, proving total vanishing.
That condition fails in the fixed-word microlocal ray.
Neither this observation nor Theorem~\ref{thm:word-quotient} proves that sewing or dynamics preserves $J_W$.

\section{Selected Abel limits: a negative and a positive example}
\label{sec:abel}

For $B_n=e^{\alpha n}$ put $r=e^{-\eps}\in(0,1)$.
The geometric family is
\begin{equation}
 q_r(n)=(1-r)r^n,\qquad
 A_r(x)=(1-r)\sum_{n\geq0}r^n x_n.
 \label{eq:abel}
\end{equation}
Every fixed law is moment determinate and the family escapes as $r\uparrow1$.
Its cluster states are more restricted than all scalar-corona states.

\begin{proposition}[Shift invariance and parity]
\label{prop:abel-shift}
For $(Sx)_n=x_{n+1}$,
\begin{equation}
 \norm{A_r\circ S-A_r}=2(1-r).
 \label{eq:shift-norm}
\end{equation}
Every weak-* cluster state is therefore shift invariant.
For the even indicator $e_n=\ind_{\{n\text{ even}\}}$,
\begin{equation}
 A_r(e)=\frac1{1+r}\longrightarrow\frac12.
 \label{eq:parity}
\end{equation}
No character of $\ell^\infty/c_0$ is a cluster state of the full Abel family.
\end{proposition}
\begin{proof}
The coefficient of $x_0$ in $A_r\circ S-A_r$ is $-(1-r)$; that of $x_m$, $m\geq1$, is $(1-r)^2r^{m-1}$.
Their absolute sum is $2(1-r)$, attained by $x_0=-1$, $x_m=1$ for $m\geq1$.
This proves~\eqref{eq:shift-norm} and shift invariance in the limit.
Summing even powers gives~\eqref{eq:parity}.
A character takes a projection to zero or one, since $\chi(e)^2=\chi(e)$.
It cannot give $[e]$ the value one half.
\end{proof}

The positive normalized shift-invariant functionals are Banach limits.
The parity conclusion and general almost-convergence theory are classical~\cite[Section 1]{Lorentz1948}.
The direct Abel calculation shows why choosing two ultrafilters cannot by itself prove reachability.

Now set
\begin{equation}
 N_k=2^{2^k},\quad b_k=k\bmod2,\qquad
 x_n=b_k\ (N_k\leq n<N_{k+1}),\quad x_n=0\ (n<N_0).
 \label{eq:block-witness}
\end{equation}

\begin{theorem}[Reachable long-block separation]
\label{thm:blocks}
For $\eps_k=(N_kN_{k+1})^{-1/2}$ and $r_k=e^{-\eps_k}$, one has
$A_{r_k}(x)-b_k\to0$.
There exist two reachable weak-* cluster states on the scalar corona agreeing on $c/c_0$ and giving $[x]$ the values zero and one.
\end{theorem}
\begin{proof}
The geometric probability of $[a,b)$ is exactly $r^a-r^b$.
Thus the probability of the $k$th block at the selected scale is
\begin{equation}
 M_k=e^{-\sqrt{N_k/N_{k+1}}}
       -e^{-\sqrt{N_{k+1}/N_k}}\longrightarrow1,
 \label{eq:block-mass}
\end{equation}
because $N_{k+1}/N_k\to\infty$.
Since $x=b_k$ there and $0\leq x\leq1$ everywhere,
\begin{equation}
 |A_{r_k}(x)-b_k|\leq1-M_k\longrightarrow0.
 \label{eq:block-bound}
\end{equation}
For even $k$ the witness tends to zero; for odd $k$ it tends to one.
Weak-* compactness gives a convergent subnet of each family.
Theorem~\ref{thm:tail} makes the states descend and agree on convergent tails.
They give the claimed distinct witness values.
No sequential compactness of the full nonseparable algebra is presumed.
\end{proof}

This is a concrete classical summability construction, not a new general theorem about Banach limits.
It supplies actual selected-state separation rather than merely ambient separation.
Both cluster states remain shift invariant.

\begin{corollary}[Two-sided bounds do not give a leading constant]
\label{cor:no-constant}
For $c_n=1+x_n$ and $B_n=e^{\alpha n}c_n$, one has $1\leq c_n\leq2$, but $\eps Z(\alpha+\eps)$ has subsequential limits one and two.
\end{corollary}
\begin{proof}
Exactly,
\begin{equation}
 \eps Z(\alpha+\eps)=\frac{\eps}{1-e^{-\eps}}
       \bigl(1+A_{e^{-\eps}}(x)\bigr).
\end{equation}
The prefactor tends to one, and Theorem~\ref{thm:blocks} gives both subsequential values.
\end{proof}

\subsection{Finite reproduction of the exact formulas}

The accompanying script evaluates exact geometric block formulas, rather than sampling an ensemble.
Table~\ref{tab:blocks} shows representative rows.
The complement bound $1-M_k$ is evaluated with $\operatorname{expm1}$ to avoid cancellation near one.
The code also checks the parity formula and fixed-word ray.

\begin{table}[htbp]
\centering
\caption{Finite exact-formula evaluations of the long-block example.
The error column bounds $|A_{r_k}(x)-b_k|$ analytically; the decimals are illustrative floating-point values.}
\label{tab:blocks}
\begin{tabular}{rrrrr}
\toprule
$k$ & $b_k$ & $\eps_k$ & $A_{r_k}(x)$ & $1-M_k$\\
\midrule
3 & 1 & $2.44140625\,10^{-4}$ & $0.9423354950$ & $0.06058704972$\\
4 & 0 & $5.96046448\,10^{-8}$ & $0.003884087113$ & $0.003898630530$\\
5 & 1 & $3.55271368\,10^{-15}$ & $0.9999847416$ & $1.525867265\,10^{-5}$\\
6 & 0 & $1.26217745\,10^{-29}$ & $2.328306437\,10^{-10}$ & $2.328306436\,10^{-10}$\\
\bottomrule
\end{tabular}
\end{table}

Last rounded digits may differ because witness and bound use different floating-point expressions.
The proof is~\eqref{eq:block-bound}, not an inference from finite rows.
Late saved witness values round to endpoints while the stable complement bound remains positive.
No exact finite endpoint equality or physical admissibility is inferred numerically.

\section{What a selected boundary-state conclusion requires}

For a unital observable algebra $\A_\partial$, define the reachable set
\begin{equation}
 \K_{\mathrm{reg}}=
 \{\Omega:\omega_{t_i}\xrightarrow{w^*}\Omega
       \text{ for a net }t_i\downarrow t_c\}.
\end{equation}
Let $\K_{\mathrm{adm}}$ impose independently stated admissibility conditions and put
$\K_{\mathrm{sel}}=\K_{\mathrm{reg}}\cap\K_{\mathrm{adm}}$.
For a unital data subalgebra $\mathcal R\subseteq\A_\partial$, ambiguity relative to those data is noninjectivity of
\begin{equation}
 \Res_{\mathcal R}:\K_{\mathrm{sel}}\to\mathsf S(\mathcal R),
 \qquad \Omega\mapsto\Omega|_{\mathcal R}.
\end{equation}
Distinct states are separated by an algebra element, but an application may require a gauge-invariant, sewing-compatible or operationally measurable witness.
Those requirements restrict the state and witness sets; scalar positivity does not supply them.

Theorem~\ref{thm:blocks} gives a noninjective selected restriction map when the scalar corona is the ambient algebra, no extra admissibility is imposed, and $\mathcal R=c/c_0$.
Parity cannot distinguish those selected states because its boundary value is fixed.
If the data algebra contains $[x]$, the constructed pair is distinguished by the data.
The conclusion therefore depends on what is held fixed.

Standard positive Hahn--Banach extension extends states from a unital $C^*$-subalgebra to its ambient unital algebra~\cite{Murphy1990}.
It preserves a difference on that subalgebra, but arbitrary extensions need not agree on other data.
If a unital positive conditional expectation $E:\A_\partial\to\mathcal Z$ exists, then $\Omega_\phi=\phi\circ E$ is a state.
Scalar states agreeing on $E(\mathcal R)$ yield ambient states agreeing on $\mathcal R$ and retain differences on witnesses in $\mathcal Z$ because $E(z)=z$.
This is conditional; the expectation is extra structure.

Only bounded weak-* closed identities automatically pass to cluster states.
A fixed identity $\omega(a)=\omega(b)$ passes by evaluation continuity, and positivity of a fixed finite matrix passes under entrywise limits.
Unbounded dynamical identities require common domains and estimates beyond weak-* compactness.
The state and operation quotients must also agree; Section~\ref{sec:word} shows why a silent change of ideal is consequential.

For an application to Virasoro topological quantum field theory, the density~\eqref{eq:density} is a source input, but actual weighted critical asymptotics, a selected packet algebra, ideal-preserving sewing and reflection maps, common dynamical domains, a physical witness and two admissible reachable states on a declared data fiber have not been constructed here.
The fixed-topology framework~\cite{CollierEberhardtZhang2023} does not prove those extra statements.
The analysis supplies tests and counterexamples, without asserting a physical gravitational completion.

\section{Concluding interpretation}

Finite exponential moments determine each regulated scalar clock.
At the actual convergence abscissa, the denominator decides whether probability remains ordinary or escapes.
Escape alone diverges every positive complexity moment; regular residual profiles control rates and, with a residual limit, give an exponential scaling law.

Coronas provide a precise home for escaping bounded states, but ambient multiplicity differs from regulator reachability.
Abel averaging gives both a uniquely averaged parity observable and long-block observables separating reachable states.
The word/total distinction requires a precise ideal: the further word quotient repairs the scalar-state construction without replacing proofs of operation preservation.
These are methods and explicit examples built from classical tools on a definite mathematical question.
They neither promote the ingredients to new general theorems nor interpret every corona state as a physical completion.

\section*{Reproducibility and authorship disclosure}

The package contains source, bibliography, standard-library Python code and saved formula evaluations.
Default reproduction needs no network, historical project path or provider service.
The central results are analytic proofs; finite output illustrates the block bounds and does not certify physical premises.
Human author and affiliation fields are intentionally blank, including PDF metadata.
Earlier work was Fable 5-led; from July 2026 onward, results were predominantly produced by GPT models, as reported by the operator.
DeepSeek and multiple GPT models also participated in review and preparation.
Human expert review has not occurred.
The operator states: ``I genuinely do not know how authorship, intellectual credit, or ownership should be attributed for these results.''
No publication-priority or journal-acceptance claim is made.
\bibliographystyle{plainnat}
\bibliography{references}
\end{document}
